\section{Zum Begriff des Nationalsozialismus}

\subsection{Die historische Gleichsetzung und ihr Fehler}

Der \emph{Nationalsozialismus} wird im öffentlichen Diskurs mit Hitler und dem Holocaust gleichgesetzt. Das ist historisch ungenau und stellt einen zentralen Bestandteil der Begriffsverdrehung dar.

Der Nationalsozialismus muss in zwei Phasen unterschieden werden:

\begin{table}[h]
\centering
\begin{tabular}{|p{0.2\textwidth}|p{0.35\textwidth}|p{0.35\textwidth}|}
\hline
\textbf{Phase} & \textbf{Früher Nationalsozialismus (1919--1933)} & \textbf{Später Nationalsozialismus (ab 1933)} \\
\hline
\textbf{Charakter} & Linke, nationalistische Arbeiterbewegung & Vom Kapital gekaufte und umgedrehte Bewegung \\
\hline
\textbf{Programm} & Sozialistische Forderungen (Verstaatlichung, Bodenreform, Gewinnbeteiligung) & Imperialismus, Krieg, Vernichtung \\
\hline
\textbf{Führung} & Linker Flügel (Strasser, Röhm) dominant & Hitler als Fassade des Kapitals \\
\hline
\textbf{Verhältnis zum Kapital} & Feindlich (Zinsknechtschaft brechen, Trusts verstaatlichen) & Werkzeug des Kapitals \\
\hline
\textbf{Ergebnis} & Bedrohung für das Kapital & Faschismus (Kapitalsteuerung) \\
\hline
\end{tabular}
\caption{Die zwei Phasen des Nationalsozialismus}
\end{table}

\subsection{Die Ursprünge: Eine linke Arbeiterbewegung}

Die NSDAP wurde 1919 als \emph{Deutsche Arbeiterpartei (DAP)} gegründet – von Arbeitern, nicht von Industriellen. Das 25-Punkte-Programm von 1920 enthielt klare sozialistische Forderungen:

\begin{itemize}
    \item Brechung der Zinsknechtschaft
    \item Verstaatlichung von Trusts
    \item Gewinnbeteiligung an Großbetrieben
    \item Bodenreform und unentgeltliche Enteignung von Boden
\end{itemize}

Die Partei hatte einen starken \emph{linken Flügel} um die Brüder Gregor und Otto Strasser, der diese sozialistischen Programmpunkte ernst nahm. Dieser Flügel wurde von Hitler systematisch zerschlagen:

\begin{itemize}
    \item Otto Strasser verließ die Partei 1930.
    \item Gregor Strasser wurde 1934 im Röhm-Putsch ermordet.
    \item Ernst Röhm, der Führer der SA, wurde ebenfalls 1934 ermordet – weil er eine soziale Revolution forderte.
\end{itemize}

\subsection{Die Zerschlagung des linken Flügels}

Die entscheidende Einsicht ist: Der Nationalsozialismus war in seinen Ursprüngen eine \emph{linke, nationalistische Arbeiterbewegung}. Der Judenhass war \emph{nicht} der Ursprung – er war Hitlers \emph{persönliche Obsession}, die er in die Partei hineintrug, um sie von innen umzuformen.

\begin{itemize}
    \item \textbf{Hitler} machte den Judenhass zum Programm, weil er Führer wurde – nicht weil die Bewegung ihn wollte.
    \item \textbf{Der linke Flügel} (Strasser, Röhm) war nicht primär antisemitisch – er war sozialistisch und nationalistisch.
    \item \textbf{Die Zerschlagung des linken Flügels} war der entscheidende Schritt, um die Bewegung für das Kapital nutzbar zu machen.
\end{itemize}

Erst \emph{nach} der Zerschlagung des linken Flügels und \emph{nach} der Machtübernahme wurde der Nationalsozialismus zum Werkzeug des Kapitals – und damit zum Faschismus im strukturellen Sinne.

\subsection{Die Umkehr: Vom Sozialismus zum Faschismus}

Diese historische Dialektik ist entscheidend:

\begin{quote}
\emph{Der Nationalsozialismus ist nicht von Anfang an faschistisch – er wird es, indem er vom Kapital vereinnahmt wird.}
\end{quote}

Der Mechanismus der Umkehr:

\begin{enumerate}
    \item \textbf{Die sozialistische Massenbewegung} bedroht das Kapital.
    \item \textbf{Das Kapital kauft die Führung} der Bewegung (Parteispenden, Medien, Einfluss).
    \item \textbf{Der linke Flügel wird zerschlagen} – die sozialistischen Ideale werden verraten.
    \item \textbf{Die Bewegung wird umgedreht} – sie wird imperialistisch, nationalistisch, gewalttätig.
    \item \textbf{Das Kapital kontrolliert die Politik} – der Faschismus ist geboren.
\end{enumerate}

Der Judenhass war \emph{nicht} der Grund für diese Umkehr – er war das \emph{Werkzeug}, das Hitler nutzte, um die Massen nach dem Verrat an den sozialistischen Idealen bei der Stange zu halten.

\subsection{Die Korrektur einer verbreiteten Fehlinterpretation}

Der öffentliche Diskurs vermischt zwei Ebenen:

\begin{itemize}
    \item \textbf{Judenhass} ist ein uraltes, eigenständiges Phänomen – nicht der Faschismus.
    \item \textbf{Hitler} machte den Judenhass zu seinem persönlichen Programm – aber er war \emph{nicht} der Grund für den Faschismus.
    \item \textbf{Der Nationalsozialismus} wurde durch den \emph{Aufkauf durch das Kapital} zum Faschismus – nicht durch Judenhass.
\end{itemize}

Die entscheidende Frage lautet daher nicht: \emph{"War der Nationalsozialismus antisemitisch?"} – das war er, weil Hitler es so wollte.

Die entscheidende Frage lautet: \emph{"Wie wurde aus einer sozialistischen Bewegung ein Werkzeug des Kapitals?"}

Die Antwort: Durch \textbf{Korruption und Aufkauf}.

\subsection{Die Parallele zur AfD}

Die heutige AfD weist strukturelle Parallelen zur frühen NSDAP auf:

\begin{itemize}
    \item Sie hat einen \emph{linken, nationalistischen Flügel} (um Björn Höcke, den sogenannten \emph{Flügel}), der soziale Forderungen stellt und sich auf die Strasser-Brüder beruft.
    \item Sie hat einen \emph{kapitalnahen Flügel} (um Alice Weidel und Tino Chrupalla), der Steuersenkungen für Reiche, Deregulierung und Abbau des Sozialstaats fordert.
    \item Sie wird zunehmend von \emph{kapitalistischen Interessen} finanziert und unterstützt – namentlich von US-amerikanischen Tech-Milliardären wie Elon Musk, Peter Thiel und dem politischen Netzwerk um Donald Trump.
\end{itemize}

Die AfD ist \emph{nicht zwingend faschistisch} – im strukturellen Sinne der Kapitalsteuerung. Aber sie \emph{wird es werden}, weil ihr kapitalnaher Flügel den linken Flügel verdrängt und sie zunehmend zum Werkzeug des Kapitals wird. Die Parallele zur NSDAP ist eindeutig: Auch sie begann als linke, nationalistische Arbeiterbewegung – und wurde vom Kapital vereinnahmt.

\subsection{Der entscheidende Unterschied zur CDU/CSU}

Der entscheidende Unterschied zur CDU/CSU ist:

\begin{itemize}
    \item \textbf{Die CDU/CSU ist bereits faschistisch} – sie ist vollständig vom Kapital durchdrungen und steuert Politik im Interesse der Industrie.
    \item \textbf{Die AfD ist auf dem Weg dorthin} – aber sie ist es noch nicht vollständig.
\end{itemize}

Die CDU/CSU ist die gegenwärtige Erscheinungsform des Faschismus im Anzug – sie wird von Lobbyisten, Parteispenden und Medien gesteuert, sie ist antidemokratisch in der Wirkung, und sie verteidigt die Interessen des Kapitals gegen die Interessen der Mehrheit.

\subsection{Zusammenfassung: Die drei Phasen der NSDAP}

Die Entwicklung der NSDAP lässt sich in drei Phasen fassen:

\begin{enumerate}
    \item \textbf{Phase 1 (1919--1923):} Linke, nationalistische Arbeiterbewegung mit sozialistischen Forderungen. Bedrohung für das Kapital.
    \item \textbf{Phase 2 (1924--1933):} Hitler zerschlägt den linken Flügel, der Judenhass wird zum Programm. Die Bewegung wird für das Kapital nutzbar gemacht.
    \item \textbf{Phase 3 (ab 1933):} Die Bewegung ist vom Kapital gekauft und umgedreht – sie wird zum Faschismus. Der Holocaust ist Hitlers persönliches Projekt, industriell durchgeführt.
\end{enumerate}

Die entscheidende Erkenntnis: \emph{Nicht der Judenhass machte den Nationalsozialismus zum Faschismus – der Aufkauf durch das Kapital war es.}

\subsection{Die universelle Wiederholung des Musters}

Das gleiche Muster wiederholt sich weltweit:

\begin{itemize}
    \item \textbf{Italien:} Fasci di Combattimento – ursprünglich sozialistisch? Nein, aber sie wurden vom Kapital gekauft und umgedreht.
    \item \textbf{Deutschland:} NSDAP – ursprünglich sozialistisch, wurde vom Kapital gekauft und umgedreht.
    \item \textbf{AfD (Gegenwart):} Ursprünglich mit linkem Flügel, wird zunehmend vom Kapital gekauft und umgedreht.
\end{itemize}

Überall dasselbe Muster: \emph{Eine Massenbewegung wird vom Kapital gekauft, umgedreht und als Werkzeug gegen die Demokratie eingesetzt.}